\section{Erweiterte Analysen und KPIs}
\label{sec:erweitert}

% Umfang: ca. 3 Seiten
% Assignment-Fragen 8-10:
%   8  What is missing to provide more detailed insights into the processes?
%   9  Where could supervised / unsupervised ML techniques help?
%   10 Can you think of useful KPIs for the processes?

\subsection{Fehlende Daten}
\label{sec:fehlende_daten}
Dass 2018 die Scope-Änderung von 2017 nicht dokumentiert wurde, ist ein Beispiel dafür, dass die Log-Daten nicht alle relevanten Informationen enthalten.
2018 wurde offenbar nicht nur die Datenerhebung erkennbar hochgeschraubt, indem weitere Fakultäten angebunden wurden, sondern auch der Prozess selbst verändert.
Mehrere ähnliche Prozesse im selben Log sorgen dafür, dass das Modell flexibler werden muss, um mehr abzudecken.
Anzunehmen ist, dass auch abseits vom Jahreswechsel zu 2018 weitere Änderungen im Prozess stattgefunden haben.
Auch verbleibt Unsicherheit darüber, ob alle Traces von vornherein aufgezeichnet oder ggf. nachträglich im Rahmen einer Art Migration nachgetragen wurden, um das Log zu vergrößern.
Prozessabläufe in International ohne Permit, falscher Permit id, mit identischer Reise und Payment sind ein Indiz dafür, dass hier Informationen ausblieben.
Wichtig ist auch, eine feste Liste gültiger Prozessabschlüsse zu haben. So sind nun mehrere Traces ohne vollständige Declaration-Bearbeitung enthalten, die nur ein Saved-Event enthalten.
Weiter ist die Rollenhistorie der Akteure nicht dokumentiert und unklar, wann welche Akteure organisatorisch eingeführt worden sind oder wie die Kompetenzen verteilt sind.
So sind die Bedeutungen des Budgets und den Limits nicht klar erkennbar.
Ebenfalls sind noch die Bedeutungen und Gründe der verschiedenen Ablehnungsoperationen unklar, um effektive Schlussfolgerungen zu treffen. \par
Bestätigt ist auch noch nicht, ob und wie Evententitäten (z.B. Declaration, Permit, Payment) wirklich miteinander verknüpft sind und auch so gelebt werden oder doch es sich nur um eine sprachliche Bündelung handelt ohne nähere Abläufe.
Mit der Ergänzung dieser Daten könnten je Zeitpunkt die wahren zugrundeliegenden Prozessvarianten identifiziert werden, was das Fitness und die Präzision unterstüzen würde und bestätigte Objektkontexte berücksichtigen kann.
% Eine kritische Reflexion darüber, welche Daten oder Attribute in den Logs fehlen.
% - keine Ablehnungsgründe
% - org:resource anonymisiert (nur STAFF MEMBER / SYSTEM) -> keine Social-Network-/Workload-Analyse
% - nur Abschluss-Events, keine Start-Events -> nur Warte-, keine Bearbeitungszeiten
% - keine Reisemerkmale (Ziel, Dauer, Kategorie)
% - keine Dokumentation der Scope-Änderung 2017 -> 2018


\subsection{Potenziale des maschinellen Lernens}
\label{sec:ml}

% Überwacht: Outcome-orientiertes Predictive Process Monitoring (Ablehnung ja/nein,
%   lange Durchlaufzeit) aus Präfix + Case-Attributen \parencite{teinemaa_outcome-oriented_2019}
% Unüberwacht: Anomalieerkennung \parencite{nolle_binet_2022}, Trace-Clustering
% Optional: Prototyp von Jesco (Python)

\todo{ML-Potenziale (Jesco)}

\subsection{Key Performance Indicators}
\label{sec:kpis}

% Definition + PQL-Formel + Werte Domestic/International \parencite{dumas_fundamentals_2018,vogelgesang_celonis_2022}
% Vorschläge: Median-Durchlaufzeit, First-Time-Right-Quote, Rework-Quote,
%             Permit-Compliance, Einreichungsverzug nach Reise, Zahlungsdauer
% Dashboard-Screenshots nach figures/celonis/

\todo{KPIs mit PQL (Eziz)}
